(1) 圆柱坐标Helmholtz方程 $(\nabla^2+\omega^2/c^2)p=0$，Bessel函数满足 $J_m''+J_m'/r+(k_r^2-m^2/r^2)J_m=0$，故 $\omega^2=c^2(k_r^2+k_z^2)$。整数 $m$ 保证绕轴一周声场单值。
(2) 线性Euler方程 $-i\omega\rho\mathbf v=-\nabla p$，故 $\mathbf v=\nabla p/(i\omega\rho)$。于是 $v_r=Pk_rJ_m'(k_rr)e^{i\cdots}/(i\omega\rho)$，$v_\phi=mp/(\omega\rho r)$，$v_z=k_zp/(\omega\rho)$，其中径向分量与声压正交相位。
(3) 平均声强 $\mathbf I=\frac12\Re(p\mathbf v^*)$，故 $I_\phi=m|p|^2/(2\omega\rho r)$，$I_z=k_z|p|^2/(2\omega\rho)$。角向流符号由 $m$ 决定；相位绕轴累积 $2\pi m$，代表涡旋绕数而非声源实际绕轴转动。
(4) $J_m(k_rr)\sim (k_rr/2)^{|m|}/|m|!$，所以 $m\ne0$ 时 $I_\phi\propto r^{2|m|-1}\to0$，轴心无发散能流；$m=0$ 时 $I_\phi=0$。无限Bessel束总能量发散，实验中须有限孔径截断。
